\documentclass[11pt]{article}
\usepackage{amsmath,amssymb,amsthm}
\newtheorem{theorem}{Theorem}
\begin{document}
\title{An eight-pole positive pool whose sign cube covers four queries}
\author{}\date{}\maketitle
\paragraph{Status.}
This is a complete elementary written construction, with exact rational
and guarded Laurent certificates. It is not Lean-checked. It concerns a
fixed finite size and does not prove an asymptotic compiler or either
accepted bound on \texttt{worstCaseHeadComplexity}.

\begin{theorem}
Let $0<t_0<t_1<t_2<t_3$, and let $0<8\epsilon<t_0$.
There are positive numbers $v_1,\ldots,v_8$ such that, for every
$(a,b,c,d)\in\{-1,1\}^4$, the function
\[
 F_{a,b,c,d}(t)=\sum_{h=1}^8\frac{v_h\tau_h}{t-h\epsilon},
 \qquad \tau=(a,b,c,d,-c,c,-a,a),
\]
has strict sign word $(d,-a,b,c)$ at $(t_0,t_1,t_2,t_3)$.
Thus the positive weighted sign cube on eight exterior poles meets
all sixteen orthants on four fixed queries. The same pool works for
all sixteen words.
\end{theorem}
\begin{proof}
Write
\[
 R(t)=(t-t_0)(t-t_2),\qquad
 P_A(t)=(t-\epsilon)(t-7\epsilon)(t-8\epsilon),\qquad
 P_C(t)=(t-3\epsilon)(t-5\epsilon)(t-6\epsilon).
\]
First choose positive single-pole functions
\[
 B(t)=\frac{v_2}{t-2\epsilon},\qquad
 D(t)=\frac{v_4}{t-4\epsilon}
\]
so that $D(t_0)-B(t_0)=1$ and $B(t_2)-D(t_2)=1$.
For explicit formulas put
\[
 q_i=(t_i-2\epsilon)(t_i-4\epsilon)\quad(i=0,2),\qquad
 L=\frac{q_0+q_2}{t_2-t_0},\qquad
 v_4=\frac{q_0+L(t_0-4\epsilon)}{2\epsilon},\qquad v_2=v_4+L.
\]
These quantities are positive and direct substitution proves both
identities. Set $S=B+D+1$, which is positive at all queries.

The ratio $k=P_C/P_A$ is strictly decreasing on $(8\epsilon,\infty)$.
Indeed, if $f(x)=1/(t-\epsilon x)$, then
\[
 \frac{d}{dt}\log k(t)
   =f(3)+f(5)+f(6)-f(1)-f(7)-f(8)<0.
\]
The last inequality follows because $f$ is strictly increasing and
strictly convex: $f(7)-f(5)>f(3)-f(1)$, and $f(8)>f(6)$.
In particular $k(t_1)>k(t_3)$. Put
\[
 u=-\frac{P_C(t_1)S(t_1)}{R(t_1)}>0,\qquad
 w=\frac{P_C(t_3)S(t_3)}{R(t_3)}>0,\qquad
 \alpha=\frac{u+w}{k(t_1)-k(t_3)}>0,\qquad
 \gamma=k(t_3)\alpha+w>0.
\]
Then
\[
 A(t)=\frac{\alpha R(t)}{P_A(t)},\qquad
 C(t)=\frac{\gamma R(t)}{P_C(t)}
\]
satisfy $C(t_i)-A(t_i)=S(t_i)$ for $i=1,3$.
Both vanish at $t_0,t_2$. Both are negative at $t_1$ and positive at
$t_3$. Consequently
\[
 -A(t_1)=|C(t_1)|+B(t_1)+D(t_1)+1,
 \qquad
 C(t_3)=A(t_3)+B(t_3)+D(t_3)+1.
\]
These are strict dominance inequalities. Together with the defining
identities for $B,D$, they show that
$aA+bB+cC+dD$ has the claimed four signs, for every choice of bits.

Finally, the three residues of $A$ at its ordered poles have signs
$+,-,+$. In fact, the residue at $p$ is $\alpha R(p)/P_A'(p)$,
where $R(p)>0$ since $p<t_0$. The derivative signs of the monic
cubic at its ordered roots are $+,-,+$. The same argument applies
to $C$. Taking their positive absolute residues as
$(v_1,v_7,v_8)$ and $(v_3,v_5,v_6)$, respectively, gives exactly
the sign assignment in the theorem. All eight residues are strictly
positive.
\end{proof}

\paragraph{Exact verification.}
\texttt{rectangular-cauchy-eight-pole-pool-check.py} verifies a saved
rational instance at queries $(77/80,39/40,79/80,1)$ and
$\epsilon=10^{-8}$. It independently checks a rational-function
continuation defined by eight score equations, using guarded exact
power-series arithmetic. Every one of the sixteen grouped words has
four positive signed leading scores, and all eight residues have
positive leading terms. The companion JSON contains the exact finite
residues, equations, and all leading coefficients. These certificates
are supplementary to the elementary proof above.
\texttt{rectangular-cauchy-eight-pole-formula-check.py} checks the
explicit formulas directly at several rational query grids and pole
scales, including all sixteen words each time.

\paragraph{Scope of the potential recursion.}
The code in the theorem consists of four disjoint pole groups of
sizes $3,1,3,1$. Each new three-pole group carries the quadratic
factor vanishing at the two earlier queries. Repeating this literal
construction by adding a pair of queries would require two new
filters vanishing at all earlier queries. Their numerator degree
and pole support would grow with the number of earlier queries;
the resulting count is quadratic. The theorem therefore supplies
no linear-pole induction.

A superconstant compiler gain would require a new way to share these
filters or use nonlinear sign preimages while preserving one fixed
pool and the required joint color counts. Neither the sixteen-word
coverage nor the large number of available sign assignments establishes
that further statement. The theorem changes no model definition,
input dimension, or atom admissibility condition.
\end{document}
